\documentclass[12pt]{article}
\usepackage[utf8]{inputenc}
\usepackage{amsmath}
\usepackage{amssymb}
\usepackage{graphicx}
\usepackage{booktabs}
\usepackage{listings}
\usepackage{xcolor}
\usepackage{hyperref}
\usepackage{geometry}
\geometry{margin=1in}

\lstset{
    language=Python,
    basicstyle=\ttfamily\small,
    keywordstyle=\color{blue}\bfseries,
    commentstyle=\color{green!60!black},
    stringstyle=\color{red},
    numbers=left,
    numberstyle=\tiny\color{gray},
    stepnumber=1,
    numbersep=5pt,
    backgroundcolor=\color{gray!10},
    frame=single,
    breaklines=true,
    breakatwhitespace=false,
    tabsize=4,
    showstringspaces=false
}

\title{Right-of-Way Costs Calculation Documentation}
\author{}
\date{}

\begin{document}

\maketitle

\section{Introduction}

The \texttt{row\_costs.py} script calculates the right-of-way (ROW) costs for a transmission line project. It computes acquisition costs (one-time purchase), holding costs (annual costs during delay period), and rental costs (annual costs during operation), then calculates both regulatory (AFUDC) and societal (present value) perspectives.

\section{Method Overview}

\subsection{Purpose}

This script calculates ROW costs by:
\begin{enumerate}
    \item Determining ROW width based on project category
    \item Calculating acres for each terrain zone
    \item Computing acquisition, holding, and rental costs per zone
    \item Aggregating costs across all zones
    \item Applying AFUDC capitalization for regulatory perspective (acquisition only)
    \item Computing present values for societal perspective
\end{enumerate}

\subsection{Calculation Approach}

The ROW cost calculation follows this structure:

\begin{align}
\text{Zone Acres} &= \frac{\text{Zone Miles} \times 5280 \times \text{ROW Width (feet)}}{43560} \\
\text{Acquisition Cost} &= \sum_{\text{zones}} \text{Zone Acres} \times \text{Acquisition Cost/Acre} \\
\text{Holding Cost} &= \sum_{\text{zones}} \text{Zone Acres} \times \text{Holding Cost/Acre/Year} \times \text{Delay Years} \\
\text{Rent Cost} &= \sum_{\text{zones}} \text{Zone Acres} \times \text{Rent Cost/Acre/Year} \times \text{Operation Years} \\
\text{Total Nominal} &= \text{Acquisition} + \text{Holding} + \text{Rent}
\end{align}

For reconductoring projects, acquisition and holding costs are zero since existing ROW is used. Only rental costs apply.

\subsection{Inputs and Outputs}

\textbf{Inputs:}
\begin{itemize}
    \item Project technical details (construction type, AC/DC, capacity, etc.)
    \item Physical details (terrain distribution by zone)
    \item ROW width for project category
    \item ROW cost parameters (acquisition, holding, rent per acre) by zone
    \item Financing parameters (WACC, inflation, AFUDC rate)
    \item Timeline (delay years, construction years, project lifetime)
\end{itemize}

\textbf{Outputs:}
\begin{itemize}
    \item Nominal costs (acquisition, holding, rent)
    \item AFUDC capitalized costs (acquisition only, regulatory perspective)
    \item Present value costs (societal perspective)
    \item Total acres required
\end{itemize}

\subsection{Integration with Other Scripts}

This script integrates with:
\begin{itemize}
    \item \texttt{yaml\_loaders.py} - Loads project configuration and ROW details (see \texttt{utilities\_documentation.tex})
    \item \texttt{financial\_utils.py} - Financial calculations (see \texttt{utilities\_documentation.tex})
    \item \texttt{csv\_output\_manager.py} - Writes results to CSV (see \texttt{csv\_output\_manager\_documentation.tex})
\end{itemize}

\section{Key Functions}

\subsection{\texttt{calculate\_zone\_costs(row\_width\_feet)}}

Calculates ROW costs for each zone where the transmission line passes.

\textbf{Parameters:}
\begin{itemize}
    \item \texttt{row\_width\_feet} (float): Right-of-way width in feet for this project category
\end{itemize}

\textbf{Returns:}
\begin{itemize}
    \item \texttt{yearly\_holding\_cost} (float): Total annual holding cost across all zones
    \item \texttt{acquisition\_cost} (float): Total one-time acquisition cost
    \item \texttt{yearly\_rent\_cost} (float): Total annual rental cost across all zones
    \item \texttt{total\_acres} (float): Total acres required across all zones
\end{itemize}

\textbf{Key Logic:}
\begin{lstlisting}
# For each zone with miles > 0:
for zone_details in row_details["right_of_way"].values():
    if zone_details["miles"] > 0:
        # Calculate acres: (miles * 5280 ft/mile * width_ft) / 43560 ft^2/acre
        zone_acres = (zone_details["miles"] * 5280 * row_width_feet) / 43560
        
        # Accumulate costs
        acquisition_cost += zone_acres * zone_details["acquisition_cost"]
        yearly_rent_cost += zone_acres * zone_details["rent_cost"]
        yearly_holding_cost += zone_acres * zone_details["hold_cost"]
\end{lstlisting}

\subsection{\texttt{main()}}

Main execution function that orchestrates the ROW cost calculation and output.

\textbf{Workflow:}
\begin{enumerate}
    \item Loads project technical details and constructs category identifier
    \item Loads ROW width for the project category
    \item Calls \texttt{calculate\_zone\_costs()} to compute costs per zone
    \item Handles reconductoring vs. new build logic
    \item Calculates AFUDC capitalized costs (acquisition only, if enabled)
    \item Calculates present values for all cost components
    \item Displays results and writes to CSV
\end{enumerate}

\textbf{Reconductoring Logic:}
\begin{lstlisting}
if reconductoring:
    # No acquisition or holding costs (existing ROW)
    total_holding_cost = 0
    acquisition_cost = 0
    # Only rent costs apply
    total_rent_cost = yearly_rent_cost * (project_lifetime + construction_years)
else:
    # New build: all costs apply
    total_holding_cost = yearly_holding_cost * delay_year
    total_rent_cost = yearly_rent_cost * (project_lifetime + construction_years)
\end{lstlisting}

\section{Example}

The following example demonstrates a typical usage of the ROW cost calculation:

\begin{lstlisting}
# Example: 500 MW Overhead AC line
# Category: "Overhead/AC/500MW/Advanced Aluminum/NA"
# 
# Inputs (from YAML):
#   - ROW width: 200 feet
#   - Terrain distribution:
#     * Zone 1 (Forested): 40 miles, $5000/acre acquisition, $50/acre/year rent
#     * Zone 2 (Farmland): 60 miles, $3000/acre acquisition, $30/acre/year rent
#   - Delay period: 2 years
#   - Project lifetime: 50 years
#   - Construction: 2 years
#
# Calculation:
#   Zone 1 acres = (40 * 5280 * 200) / 43560 = 969.7 acres
#   Zone 2 acres = (60 * 5280 * 200) / 43560 = 1454.5 acres
#   Total acres = 2424.2 acres
#
#   Acquisition = (969.7 * $5000) + (1454.5 * $3000) = $7.3M
#   Holding (2 years) = (969.7 * $50 * 2) + (1454.5 * $30 * 2) = $0.2M
#   Rent (52 years) = (969.7 * $50 * 52) + (1454.5 * $30 * 52) = $5.2M
#   Total nominal = $12.7M
#
# Output:
#   - Nominal cost: $12.7M
#   - AFUDC capitalized (acquisition): $7.5M (if AFUDC enabled)
#   - Present value: $11.2M (discounted at real WACC)
\end{lstlisting}

\textbf{Integration Example:}

\begin{lstlisting}
# When called from ctcc.py:
csv_manager = CTCCOutputManager()
results = {
    "total_nominal": total_nominal_cost,
    "total_afudc": acquisition_capitalized + total_holding_cost + total_rent_cost,
    "total_pv": total_pv_cost,
    "acquisition_nominal": acquisition_cost,
    "holding_nominal": total_holding_cost,
    "rent_nominal": total_rent_cost,
}
csv_manager.add_row_costs(results)
csv_manager.write_batch_summary()
\end{lstlisting}

\section{Key Implementation Details}

\subsection{AFUDC Eligibility}

Only acquisition costs are AFUDC-eligible (capitalized to plant cost). Holding costs are operating expenses during the delay period, and rent costs are operational period expenses. Neither holding nor rent costs receive AFUDC treatment.

\subsection{Reconductoring Handling}

For reconductoring projects, acquisition and holding costs are set to zero since existing ROW is used. Only rental costs apply during the operational period.

\subsection{Present Value Timing}

\begin{itemize}
    \item \textbf{Holding costs}: Incurred annually during delay period (years 1 to delay\_year)
    \item \textbf{Acquisition costs}: One-time payment at end of delay period (year delay\_year)
    \item \textbf{Rent costs}: Incurred annually during operation (years delay\_year+1 to delay\_year+construction\_years+project\_lifetime)
\end{itemize}

\subsection{Acres Calculation}

The script converts miles and ROW width to acres using the standard conversion:
\[
\text{Acres} = \frac{\text{Miles} \times 5280 \text{ ft/mile} \times \text{ROW Width (ft)}}{43560 \text{ ft}^2/\text{acre}}
\]

\end{document}

